
\subsubsection{2.1 Dataset Versioning Infrastructure}

Data Version Control (DVC) \citep{Kuprieievetal2020} provides Git-like versioning for datasets, assigning unique hashes to each version. While DVC tracks changes, it does not classify their severity—a metadata correction and a complete data redistribution receive identical version treatment. HuggingFace Datasets \citep{Lhoestetal2021} uses Git commit identifiers for version tracking, enabling reproducibility through version pinning but providing no automated severity classification.

\subsubsection{2.2 Distribution Shift Detection}

The Kolmogorov-Smirnov test \citep{Massey1951} measures the maximum distance between empirical cumulative distribution functions of two samples. Maximum Mean Discrepancy \citep{Grettonetal2012} measures distribution distance in a reproducing kernel Hilbert space. Rabanser et al. [2019] benchmarked drift detection methods and found that no single metric performs optimally across all shift types, noting that threshold selection depends on application context.

González-Cebrián et al. [2024] proposed using PCA and autoencoders for dataset version detection but did not classify severity levels. TorchDrift [Schröder et al., 2021] implements KS and MMD tests with p-value computation but does not map drift magnitudes to semantic version labels.

\subsubsection{2.3 Dataset Evolution Effects}

Recht et al. [2019] documented that ImageNet classifiers experienced 10-15\% accuracy reduction on ImageNet-v2, a carefully reproduced test set. This demonstrated that dataset curation differences produce measurable performance impacts. The CIFAR-10.1 dataset \citep{Rechtetal2018} similarly showed performance degradation from distribution shift.

\subsubsection{2.4 Positioning}

This work is the first to test automated MAJOR/MINOR/PATCH classification for datasets using statistical drift with fixed thresholds. Unlike prior work that detects drift or tracks versions, this work attempted to classify severity using predefined thresholds derived from ImageNet literature. The 28.57\% accuracy result provides empirical evidence that this approach does not generalize.
